\documentclass[../../../include/open-logic-section]{subfiles}

\begin{document}
\olfileid[id]{sth}{ordinals}{zfm}
\olsection{$\ZFminus$: suatu tonggak penting}

Pertanyaan tentang bagaimana Penggantian dapat dibenarkan (jika memang
dapat) tidaklah sederhana. Karena itu, kita akan menundanya hingga
\olref[replacement][]{chap}. Namun, dengan penambahan Penggantian, kita telah
mencapai tonggak penting lain. Kini kita mempunyai semua aksioma yang
diperlukan bagi teori $\ZFminus$. Secara terperinci:

\begin{defn}
Teori $\ZFminus$ mempunyai aksioma-aksioma berikut: Ekstensionalitas,
Gabungan, Pasangan, Himpunan Kuasa, Ketakhinggaan, serta semua instans
skema Pemisahan dan Penggantian. Dengan kata lain, $\ZFminus$ menambahkan
Penggantian pada $\Zminus$.
\end{defn}
\noindent
Nama ini merupakan singkatan bagi teori himpunan
\emph{Zermelo--Fraenkel} (\emph{dikurangi} sesuatu yang akan kita bahas
kemudian). Fraenkel mendapat kehormatan tersebut karena ia dianggap berjasa
merumuskan Penggantian dalam \citeyear{Fraenkel1922}, meskipun perumusan
tepat yang pertama berasal dari \citet{Skolem1922}.

\end{document}
